%% kapa: exam=1 round=19 year=2008 period=2 form=A subject=economics subject_ko=경제원론 seq=2 number=2
%% source: Q-Net 감정평가사 기출문제 (cst00309) artlSeq=1001962 file=extracted/1008/1001962/2008년 제19회 감정평가사 1차시험 기출문제/2교시.hwp
%% conversion: kordoc 4.24.51 -> tools/gen1.py (auto); review=proofread-sonnet; flags=none
\begin{question}{2}
구매력평가설과 이자율평가설이 성립한다고 하자. 한국의 명목이자율은 연 4\%, 미국의 명목이자율은 연 2\%이며, 한국의 물가상승률은 연 5\%로 예상될 경우에 옳은 것을 모두 고르면?
\begin{framedtext}
ㄱ. 한국의 실질이자율은 연 -1\%이다. \par
ㄴ. 한국과 미국의 실질이자율은 동일하다. \par
ㄷ. 1년 후 원/달러 환율은 2\% 상승할 것으로 예상된다. \par
ㄹ. 미국의 물가상승률은 연 2\%일 것으로 예상된다. \par
\end{framedtext}
\begin{choices}
  \item ㄱ, ㄴ
  \item ㄱ, ㄷ
  \item ㄴ, ㄹ
  \item ㄱ, ㄴ, ㄷ
  \item ㄴ, ㄷ, ㄹ
\end{choices}
\end{question}
\answer{4}
